\subsection{Otras variantes lineales conservadas}

En \texttt{ya emp/OLS.py} se emplea \texttt{datos\_ventanas.csv}. Se elimina la fecha objetivo de las entradas, se ajusta un único modelo con el primer 70\,\% de las filas y se evalúa el 30\,\% restante. El escalador se estima exclusivamente con entrenamiento. El generador \texttt{gen\_csv.py} construye filas con 24 registros de historia, cada uno con cinco variables de calendario y la generación, y toma como objetivo el registro siguiente.

En \texttt{dificil/ols.py} y \texttt{dificil/ridge\_dif.py} se usa \texttt{datos\_backtesting.csv}. Ambos ajustan una sola vez sobre las primeras $48-24=24$ filas y recorren bloques solapados de 72 filas, con avance de una. Ridge fija $\alpha=1$. Se excluyen los identificadores, la fecha objetivo y la respuesta de las entradas, y se estandariza con entrenamiento. La actualización modifica los rezagos de generación de la siguiente fila usando los valores que esa fila ya contiene y la predicción actual; no arrastra explícitamente todo el historial predicho. Por ello, este mecanismo no debe describirse como una recursión completa de 72 pasos sin observaciones futuras.

En \texttt{dificil/ols2.py} y \texttt{dificil/ridge\_2.py} se extrae directamente la columna de generación de \texttt{datos.csv}, sin filtrar área ni ordenar fechas en esos programas. Se reserva el primer 80\,\% para ajustar y el resto para evaluar, con 24 valores históricos y horizontes recursivos de 72 pasos. OLS completa el historial inicial con datos de entrenamiento cuando el origen está entre las primeras 23 posiciones de prueba. Ridge solo contempla ese empalme para el primer origen: en los siguientes 23, el corte del historial no garantiza las 24 entradas requeridas por el escalador. Esta limitación impide asumir que dicha versión produce todas las métricas previstas. Las etiquetas de horas de estos experimentos requieren confirmar que las filas forman una única serie horaria.
